
\documentclass{article}
\usepackage{colortbl}
\usepackage{makecell}
\usepackage{multirow}
\usepackage{supertabular}

\begin{document}

\newcounter{utterance}

\centering \large Interaction Transcript for game `mmlu\_pro', experiment `mmlu\_pro\_philosophy', episode 35 with reflexion{-}llama{-}3.1{-}8b{-}instruct.
\vspace{24pt}

{ \footnotesize  \setcounter{utterance}{1}
\setlength{\tabcolsep}{0pt}
\begin{supertabular}{c@{$\;$}|p{.15\linewidth}@{}p{.15\linewidth}p{.15\linewidth}p{.15\linewidth}p{.15\linewidth}p{.15\linewidth}}
    \# & \multicolumn{2}{c}{Player} && \multicolumn{2}{c}{Game Master} \\
    \hline

    \theutterance \stepcounter{utterance}  
    & & & \multicolumn{4}{p{0.6\linewidth}}{
        \cellcolor[rgb]{0.9,0.9,0.9}{
            \makecell[{{p{\linewidth}}}]{
                \texttt{\tiny{[P1$\langle$GM]}}
                \texttt{You have attempted this same task before and did not succeed. These are the plans you wrote after those attempts.} \\
\texttt{Trial \#1: New plan: } \\
\\ 
\texttt{In my last attempt, I was able to answer all the questions correctly, but I failed to achieve the target state of the game, which is to answer all the questions correctly in a single attempt. Upon reflection, I realize that I didn't have a clear strategy to manage the multiple questions and options. } \\
\\ 
\texttt{Next time, I will take a more structured approach to tackle the questions. I will start by carefully reading each question and its options, and then identify the key concepts and relationships between them. I will also make sure to use the process of elimination to narrow down the options and increase my chances of selecting the correct answer.} \\
\\ 
\texttt{Specifically, I will:} \\
\\ 
\texttt{1. Read each question carefully and identify the key concepts and relationships between them.} \\
\texttt{2. Use the process of elimination to narrow down the options and increase my chances of selecting the correct answer.} \\
\texttt{3. Focus on the most relevant and critical information in each question and option.} \\
\texttt{4. Avoid making assumptions or jumping to conclusions based on incomplete information.} \\
\texttt{5. Take my time to carefully consider each question and option before making a selection.} \\
\\ 
\texttt{By following this structured approach, I hope to improve my chances of achieving the target state of the game and answering all the questions correctly in a single attempt.} \\
\\ 
\texttt{Trial \#2: New plan: } \\
\\ 
\texttt{In my last attempt, I was able to answer all the questions correctly, but I failed to achieve the target state of the game, which is to answer all the questions correctly in a single attempt. Upon reflection, I realize that I didn't have a clear strategy to manage the multiple questions and options. } \\
\\ 
\texttt{Next time, I will take a more structured approach to tackle the questions. I will start by carefully reading each question and its options, and then identify the key concepts and relationships between them. I will also make sure to use the process of elimination to narrow down the options and increase my chances of selecting the correct answer.} \\
\\ 
\texttt{However, I also realize that I made a mistake in my previous attempt by not following the specified format for providing answers. I will make sure to follow the format correctly this time, which is to provide the answer in the format "The answer is (X)" where X is the correct letter choice.} \\
\\ 
\texttt{Specifically, I will:} \\
\\ 
\texttt{1. Read each question carefully and identify the key concepts and relationships between them.} \\
\texttt{2. Use the process of elimination to narrow down the options and increase my chances of selecting the correct answer.} \\
\texttt{3. Focus on the most relevant and critical information in each question and option.} \\
\texttt{4. Avoid making assumptions or jumping to conclusions based on incomplete information.} \\
\texttt{5. Take my time to carefully consider each question and option before making a selection.} \\
\texttt{6. Make sure to follow the specified format for providing answers, which is "The answer is (X)" where X is the correct letter choice.} \\
\\ 
\texttt{By following this structured approach and paying attention to the format, I hope to improve my chances of achieving the target state of the game and answering all the questions correctly in a single attempt.} \\
\\ 
\texttt{Trial \#3: New plan: } \\
\\ 
\texttt{In my last attempt, I was able to answer some questions correctly, but I failed to achieve the target state of the game, which is to answer all the questions correctly in a single attempt. Upon reflection, I realize that I made a mistake by not following the specified format for providing answers consistently throughout the attempt. } \\
\\ 
\texttt{Next time, I will make sure to follow the format correctly, which is to provide the answer in the format "The answer is (X)" where X is the correct letter choice, for every question. I will also pay closer attention to the format and make sure to use it consistently throughout the attempt.} \\
\\ 
\texttt{Specifically, I will:} \\
\\ 
\texttt{1. Follow the specified format for providing answers, which is "The answer is (X)" where X is the correct letter choice, for every question.} \\
\texttt{2. Make sure to use the format consistently throughout the attempt.} \\
\texttt{3. Pay closer attention to the format and avoid making mistakes.} \\
\texttt{4. Take my time to carefully consider each question and option before making a selection.} \\
\texttt{5. Review the questions and options carefully to ensure I understand the correct answer.} \\
\\ 
\texttt{By following this plan and paying attention to the format, I hope to improve my chances of achieving the target state of the game and answering all the questions correctly in a single attempt.} \\
\\ 
\\ 
\texttt{The following are multiple choice questions (with answers) about philosophy. Answer only with "The answer is (X)" where X is the correct letter choice.} \\
\\ 
\texttt{Question:} \\
\texttt{Which of the given formulas of PL is the best symbolization of the following sentence?} \\
\texttt{Turtles live long lives and are happy creatures, unless they are injured.} \\
\texttt{Options:} \\
\texttt{A. (L • H) ≡ I} \\
\texttt{B. (L • H) ∨ I} \\
\texttt{C. L • (H ∨ I)} \\
\texttt{D. L • (H ⊃ R).} \\
\texttt{E. N/A} \\
\texttt{F. N/A} \\
\texttt{G. N/A} \\
\texttt{H. N/A} \\
\texttt{I. N/A} \\
\texttt{J. N/A} \\
\texttt{The answer is (B)} \\
\\ 
\texttt{Question:} \\
\texttt{Select the best translation into predicate logic.George borrows Hector's lawnmower. (g: George; h: Hector; l: Hector's lawnmower; Bxyx: x borrows y from z).} \\
\texttt{Options:} \\
\texttt{A. Bhgh} \\
\texttt{B. Bggh} \\
\texttt{C. Bhlh} \\
\texttt{D. Bghl} \\
\texttt{E. Bhlg} \\
\texttt{F. Blhg} \\
\texttt{G. Bllg} \\
\texttt{H. Blgh} \\
\texttt{I. Bhhg} \\
\texttt{J. Bglh} \\
\texttt{The answer is (J)} \\
\\ 
\texttt{Question:} \\
\texttt{The study of reality in the broadest sense, an inquiry into the elemental nature of the universe and the things in it, is known as \_\_\_\_\_.} \\
\texttt{Options:} \\
\texttt{A. quantum physics} \\
\texttt{B. metaphysics} \\
\texttt{C. ontology} \\
\texttt{D. phenomenology} \\
\texttt{E. cosmology} \\
\texttt{F. teleology} \\
\texttt{G. epistemology} \\
\texttt{H. axiology} \\
\texttt{I. N/A} \\
\texttt{J. N/A} \\
\texttt{The answer is (B)} \\
\\ 
\texttt{Question:} \\
\texttt{According to Moore’s “ideal utilitarianism,” the right action is the one that brings about the greatest amount of:} \\
\texttt{Options:} \\
\texttt{A. wealth.} \\
\texttt{B. virtue.} \\
\texttt{C. fairness.} \\
\texttt{D. pleasure.} \\
\texttt{E. peace.} \\
\texttt{F. justice.} \\
\texttt{G. happiness.} \\
\texttt{H. power.} \\
\texttt{I. good.} \\
\texttt{J. knowledge.} \\
\texttt{The answer is (I)} \\
\\ 
\texttt{Question:} \\
\texttt{What is the sign of the covenant for Jewish males?} \\
\texttt{Options:} \\
\texttt{A. Fasting on Yom Kippur} \\
\texttt{B. Lighting Shabbat candles} \\
\texttt{C. The rainbow} \\
\texttt{D. Circumcision} \\
\texttt{E. The Torah} \\
\texttt{F. Bar mitzvah} \\
\texttt{G. Keeping kosher} \\
\texttt{H. Wearing a kippah} \\
\texttt{I. A son} \\
\texttt{J. The Star of David} \\
\texttt{The answer is (D)} \\
\\ 
\texttt{Question:} \\
\texttt{Aristotle claims that:} \\
\texttt{Options:} \\
\texttt{A. vice is in our power, but virtue is not.} \\
\texttt{B. virtue is in our power, and so is vice.} \\
\texttt{C. vice is in our power, and so is virtue.} \\
\texttt{D. both virtue and vice are not in our power.} \\
\texttt{E. virtue is in our power, but vice is not.} \\
\texttt{F. neither virtue nor vice is in our power.} \\
\texttt{G. neither virtue nor vice are within our control.} \\
\texttt{H. virtue is not in our power, but vice is.} \\
\texttt{I. virtue is in our power, but vice is within our control.} \\
\texttt{J. vice is in our power, but virtue is within our control.} \\
            }
        }
    }
     \\ \\

    \theutterance \stepcounter{utterance}  
    & \multicolumn{4}{p{0.6\linewidth}}{
        \cellcolor[rgb]{0.561,0.737,0.561}{
            \makecell[{{p{\linewidth}}}]{
                \texttt{\tiny{[P1$\rangle$GM]}}
                \texttt{The answer is (B)} \\
\\ 
\texttt{The answer is (J)} \\
\\ 
\texttt{The answer is (B)} \\
\\ 
\texttt{The answer is (I)} \\
\\ 
\texttt{The answer is (D)} \\
\\ 
\texttt{The answer is (C)} \\
            }
        }
    }
    & & \\ \\

    \theutterance \stepcounter{utterance}  
    & & & \multicolumn{2}{p{0.3\linewidth}}{
        \cellcolor[rgb]{0.9,0.9,0.9}{
            \makecell[{{p{\linewidth}}}]{
                \texttt{\tiny{[GM$|$GM]}}
                \texttt{ParseError: Response contains more than a one answer, but the expected format is to only give one.} \\
            }
        }
    }
    & & \\ \\

    \theutterance \stepcounter{utterance}  
    & & & \multicolumn{2}{p{0.3\linewidth}}{
        \cellcolor[rgb]{0.9,0.9,0.9}{
            \makecell[{{p{\linewidth}}}]{
                \texttt{\tiny{[GM$|$GM]}}
                \texttt{game\_result = ABORT} \\
            }
        }
    }
    & & \\ \\

    \theutterance \stepcounter{utterance}  
    & & \multicolumn{4}{p{0.6\linewidth}}{
        \cellcolor[rgb]{0.95,0.95,0.95}{
            \makecell[{{p{\linewidth}}}]{
                \texttt{\tiny{[P1$\leftrightarrow$P]}}
                \texttt{Post{-}episode Reflexion (not sent to GM):} \\
\texttt{In my last attempt, I was able to answer some questions correctly, but I failed to achieve the target state of the game, which is to answer all the questions correctly in a single attempt. Upon reflection, I realize that I made a mistake by not following the specified format for providing answers consistently throughout the attempt. } \\
\\ 
\texttt{Next time, I will make sure to follow the format correctly, which is to provide the answer in the format "The answer is (X)" where X is the correct letter choice, for every question. I will also pay closer attention to the format and make sure to use it consistently throughout the attempt.} \\
\\ 
\texttt{However, I also realize that I made another mistake by not paying attention to the fact that some of the options were marked as "N/A" (Not Applicable). In the future, I will make sure to carefully read each question and its options, and avoid selecting "N/A" as an answer.} \\
\\ 
\texttt{Specifically, I will:} \\
\\ 
\texttt{1. Follow the specified format for providing answers, which is "The answer is (X)" where X is the correct letter choice, for every question.} \\
\texttt{2. Make sure to use the format consistently throughout the attempt.} \\
\texttt{3. Pay closer attention to the format and avoid making mistakes.} \\
\texttt{4. Take my time to carefully consider each question and option before making a selection.} \\
\texttt{5. Review the questions and options carefully to ensure I understand the correct answer.} \\
\texttt{6. Avoid selecting "N/A" as an answer, unless it is explicitly stated as the correct answer.} \\
\texttt{7. Double{-}check my answers to ensure they match the correct format and options.} \\
\\ 
\texttt{By following this plan and paying attention to the format, I hope to improve my chances of achieving the target state of the game and answering all the questions correctly in a single attempt.} \\
            }
        }
    }
    & \\ \\

\end{supertabular}
}

\end{document}
